\documentclass[12pt]{article}

\usepackage{amsmath}
\usepackage{amssymb}
\usepackage{amsthm}
\newtheorem{theorem}{Theorem}[section]
\newtheorem{definition}{Definition}[section]


\begin{document}
\section{measure theory prerequisites}

\begin{definition}
	Let $S$ be a set, an \textbf{algebra} on a set of subsets of $S$: $\Sigma_0$ is called an algebra on $S$ (or an algebra on $2^S$) if 
	\begin{enumerate}
		\item $S \in \Sigma_0$ 
	\item $F \in \Sigma_0 \implies F^c := S\setminus F \in \Sigma_0$ 
	\item $F, G \in \Sigma_0 \implies F \cup G \in \Sigma_0$ 
		
	\end{enumerate}
	From these rules, obvious other rules follows such as, $\emptyset \in \Sigma_0$ obviously since $\emptyset = S^c \in \Sigma_0$. Also letting $F, G \in \Sigma_0$, then $(F^c \cup G^c)^c = F\cap G \in \Sigma_0$ 
	Thus, an algebra on S is a family of subsets of $S$ that are stable under finitely many set operations.
\end{definition}
This definition is algebraic in the sense that $\cap$ is the product operation and symmetric difference being 
\begin{align*}
	A \Delta B := (A \cup B)\setminus(A \cap B)
\end{align*}
We can extend this definition to form a $\sigma$-algebra on $S$ which takes the rules and extends them towards their countable definition. That means, if $\Sigma$ is an algebra on $S$, such that whenever $F_n \in \Sigma$ ($n \in \mathbb{N}$) then 
\begin{align*}
	\bigcup_{n} F_n \in \Sigma
\end{align*}
Obviously set intersection follows
\begin{align*}
	\bigcap_n F_n = \left( \bigcup_n F_n^c \right)^c \in \Sigma
\end{align*}
\subsection{measurable space}
A pair $(S, \Sigma)$, where $S$ is a set and $\Sigma$ is a $\sigma$-algebra on $S$, is called a measurable space. An element of $\Sigma$ is called a $\Sigma$-measurable subset of $S$. 
\begin{definition}[$\sigma$-algebra generated by a class $\mathcal{C}$ of subsets]
	Let $\mathcal{C}$ be a class of subsets of $S$, then $\sigma(\mathcal{C})$, the $\sigma$-algebra generated by $\mathcal{C}$ is the smallest $\sigma$-algebra $\Sigma$ on $S$ such that $\mathcal{C} \subseteq \Sigma$. 
	
\end{definition}
An easy example would be a borel $\sigma$-algebra, $\mathcal{B}(S)$. Indeed, letting $S$ be a topological space, $B(S)$ is the borel $\sigma$-algebra generated by the family of open subsets of $S$. $\mathcal{B}(\mathbb{R})$ is the most important of all $\sigma$-algebras, oftenhand shorthanded by $\mathcal{B}$.
Consider the collection 
\begin{align*}
	\pi(\mathbb{R}) := \{ (-\infty, x] \colon x \in \mathbb{R}\}
\end{align*}
We claim that \[\mathcal{B} = \sigma(\pi(\mathbb{R})\]
\begin{proof}
	\begin{itemize}
		\item ($\supset$) Since they are both $\sigma$-algebras, we just need to prove that they contain the same elements. We can do this exhaustively.  Taking an element of $\sigma(\pi(\mathbb{R}))$ as $(-\infty, x)$ for some $x$, we notice that this is an element of $\mathcal{B}$, since $(-\infty, x] =\bigcap_{n \in \mathbb{N}} \left(-\infty, x+\frac{1}{n}\right)$
		\item $\subset$ all that remains to be prove that is every open set $G$ of $\mathcal{B}$ is in $\sigma(\pi(\mathbb{R}))$, we recall from first problem set of analysis 2 that every open set $G$ is a countable union of open intervals, that means we just need to show that any $(a,b)$ for $a,b \in \mathbb{R}$ with $a< b$ is in $\sigma(\pi(\mathbb{R}))$. But for any $u$ with $u > a$, we have that 
			\begin{align*}
				(a,u] =(-\infty,u] \cap (a,+\infty) = (-\infty,u] \cap (-\infty, a]^c \in \sigma(\pi(\mathbb{R}))
			\end{align*}

			but for $\varepsilon = \frac{1}{2}(b-a)$, we get that 
			\begin{align*}
				(a,b) = \bigcup_n (a,b-\varepsilon \frac{1}{n}]
			\end{align*}
			to recover the open interval $(a,b)$ in $\mathbb{R}$

	\end{itemize}	
	\begin{definition}[set functions]
		Let $S$ be a set, letting 
	\end{definition}
\end{proof}



\end{document}
